\documentclass[border=10pt,tikz]{standalone}
\usepackage{fontspec}
\setmainfont{Apple SD Gothic Neo}
\setsansfont{Apple SD Gothic Neo}
\usepackage{tikz}
\usetikzlibrary{positioning,arrows.meta,shapes,fit,backgrounds}
\begin{document}
% _design.tex — shared TikZ design system for all blog diagrams
% Inputs nothing on its own — meant to be \input{} from each diagram .tex
% AFTER \begin{document}, BEFORE \begin{tikzpicture}.
%
% Provides a unified visual language across GoF / DSA / EMC++ / Beautiful series:
%   - Colors: muted, accessible, color-coded by role
%   - Node styles: consistent sizing, rounded, subtle borders
%   - Edge styles: UML-style inheritance, aggregation, plus dashed alternates
%
% Usage:
%   \documentclass[border=10pt,tikz]{standalone}
%   \usepackage{tikz}
%   \usetikzlibrary{positioning,arrows.meta,shapes,fit,backgrounds}
%   \begin{document}
%   \input{../_design.tex}     % adjust relative path
%   \begin{tikzpicture}[blog]  % apply 'blog' style
%   ...

\definecolor{absfill}{RGB}{220,232,247}   % cool blue
\definecolor{absbord}{RGB}{72,118,182}
\definecolor{confill}{RGB}{249,222,222}   % warm red
\definecolor{conbord}{RGB}{197,93,93}
\definecolor{impfill}{RGB}{223,238,222}   % soft green
\definecolor{impbord}{RGB}{99,154,93}
\definecolor{clifill}{RGB}{251,239,205}   % pale yellow
\definecolor{clibord}{RGB}{200,167,72}
\definecolor{hifill} {RGB}{253,224,200}   % accent orange
\definecolor{hibord} {RGB}{210,135,68}
\definecolor{nodefill}{RGB}{250,250,250}
\definecolor{nodebord}{RGB}{120,120,120}
\definecolor{edgec}  {RGB}{80,80,80}
\definecolor{edgesc} {RGB}{145,145,145}
\definecolor{groupbord}{RGB}{200,200,200}

\tikzset{
  blog/.style={
    font=\sffamily\small,
    every node/.style={font=\sffamily\small}
  },
  % --- node roles ---
  cls/.style={
    rectangle, rounded corners=2pt,
    draw=nodebord, line width=0.4pt,
    minimum height=8mm, minimum width=28mm,
    inner sep=3pt, fill=nodefill, align=center
  },
  abs/.style={cls, draw=absbord, fill=absfill, font=\sffamily\itshape\small},
  con/.style={cls, draw=conbord, fill=confill},
  imp/.style={cls, draw=impbord, fill=impfill},
  cli/.style={cls, draw=clibord, fill=clifill},
  hi/.style={cls, draw=hibord, fill=hifill, font=\sffamily\bfseries\small},
  % --- circle (for trees / graph nodes) ---
  vx/.style={
    circle, draw=absbord, line width=0.4pt,
    minimum size=8mm, fill=absfill, inner sep=0pt
  },
  vxc/.style={vx, draw=conbord, fill=confill},
  vxh/.style={vx, draw=hibord, fill=hifill, font=\sffamily\bfseries},
  % --- decision diamond ---
  q/.style={
    diamond, draw=clibord, line width=0.4pt, fill=clifill,
    inner sep=2pt, aspect=2.4, align=center
  },
  % --- edges ---
  inh/.style={                                  % UML inheritance
    -{Triangle[length=3mm,width=3mm,open]},
    draw=edgec, line width=0.5pt
  },
  agg/.style={                                  % UML aggregation
    {Diamond[length=2.5mm,width=2.5mm,open]}-{Stealth[length=2.2mm]},
    draw=edgec, line width=0.45pt
  },
  uses/.style={                                 % directed reference
    -{Stealth[length=2.2mm]}, draw=edgec, line width=0.5pt
  },
  arr/.style={                                  % plain arrow (graphs/trees)
    -{Stealth[length=2mm]}, draw=edgec, line width=0.45pt
  },
  edge/.style={                                 % undirected (trees)
    draw=edgec, line width=0.45pt
  },
  alt/.style={                                  % dashed: similar/alt
    -{Stealth[length=2mm]}, dashed, draw=edgesc, line width=0.45pt
  },
  ret/.style={                                  % dashed creates/returns
    -{Stealth[length=2mm]}, dashed, draw=edgesc, line width=0.45pt
  },
  thr/.style={                                  % threading link (red dashed)
    ->, dashed, draw=conbord, line width=0.5pt
  },
  % --- group / region ---
  group/.style={
    rectangle, rounded corners=4pt,
    draw=groupbord, line width=0.4pt,
    inner sep=10pt, fill=black!2
  },
  glabel/.style={font=\sffamily\bfseries\small, anchor=north west, inner sep=2pt}
}

% Korean font convenience macro — included by every Hangul diagram
\providecommand{\KOR}{}

\begin{tikzpicture}[blog,
  stg/.style={cls, minimum width=124mm, text width=118mm, align=center, font=\sffamily\small},
  leaf/.style={cls, minimum width=58mm, text width=54mm, align=center, font=\sffamily\small},
  elab/.style={font=\sffamily\scriptsize, text=mutec, anchor=west}
]

\node[font=\sffamily\bfseries\large] at (0, 2.0) {Chain of Trust: 단계별 검증 경계};

% Trust anchor
\node[stg, draw=impbord, fill=impfill, minimum height=13mm] (root) at (0, 0.6)
  {\textbf{Trust anchor}: SoC ROM의 고정 검증 로직\\[2pt]{\scriptsize 제품 정책에 따라 OTP/fuse 등 protected storage의 key 또는 key hash 추가 (chip family·lifecycle별)}};

% TF-A TBB
\node[stg, draw=absbord, fill=absfill, minimum height=21mm] (tfa) at (0, -2.0)
  {\textbf{TF-A Trusted Board Boot} {\scriptsize (\texttt{TRUSTED\_BOARD\_BOOT=1} 빌드)}\\[2pt]
   {\scriptsize BL1: BL2 content certificate의 issuer key hash를 ROTPK hash와 비교한 뒤 BL2 검증}\\[2pt]
   {\scriptsize BL2: trusted key certificate로 BL31·BL32, non-trusted key로 BL33 검증}\\[2pt]
   {\scriptsize\color{mutec} TF-A 문서의 기본 순서이며, 역할과 대상 image는 platform port에 따라 다름}};

% U-Boot FIT
\node[stg, draw=clibord, fill=clifill, minimum height=15mm] (uboot) at (0, -5.0)
  {\textbf{U-Boot (BL33) $\rightarrow$ FIT}: image hash · signature 검증\\[2pt]
   {\scriptsize control FDT의 \texttt{/signature} node에 public key, \texttt{required} 속성으로 검증 요구}\\[2pt]
   {\scriptsize\color{mutec} configuration signature와 image signature는 서로 다른 선택지}};

% Linux kernel
\node[stg, draw=neutbord, fill=neutfill, minimum height=9mm] (kern) at (0, -7.3)
  {\textbf{Linux kernel}};

% Leaves
\node[leaf, draw=prpbord, fill=prpfill, minimum height=15mm] (mod) at (-3.2, -9.4)
  {\textbf{module signing}\\[2pt]{\scriptsize enforcement가 켜진 설정에서만}\\[2pt]{\scriptsize unsigned·unknown-key module 거부}};
\node[leaf, draw=hibord, fill=hifill, font=\sffamily\small, minimum height=15mm] (vrt) at (3.2, -9.4)
  {\textbf{dm-verity (root filesystem)}\\[2pt]{\scriptsize hash tree로 읽기 내용 검증}\\[2pt]{\scriptsize root hash 전달 경로는 제품 정책}};

% Edges with conditions
\draw[arr, line width=0.8pt] (root.south) -- (tfa.north)
  node[midway, elab, xshift=2pt] {ROM이 첫 firmware의 format·인증 정책을 지원할 때};
\draw[arr, line width=0.8pt] (tfa.south) -- (uboot.north)
  node[midway, elab, xshift=2pt] {BL33으로 U-Boot 실행};
\draw[arr, line width=0.8pt] (uboot.south) -- (kern.north)
  node[midway, elab, xshift=2pt] {검증된 kernel · DTB};
\draw[arr, line width=0.8pt] (kern.south -| mod.north) -- (mod.north);
\draw[arr, line width=0.8pt] (kern.south -| vrt.north) -- (vrt.north);

% Footer note
\node[font=\sffamily\scriptsize, text=conbord, text width=124mm, align=center] at (0, -11.0)
  {각 화살표는 해당 단계가 빌드·설정되고 실패 시 실행을 막을 때만 성립합니다.\\[2pt]앞 단계가 빠지면 뒤 단계의 검증은 초기 firmware를 보호하지 못합니다.};

\end{tikzpicture}
\end{document}
